\chapter{Angiotensin Receptor Blockers (ARBs)}
\index{Angiotensin Receptor Blockers (ARBs)}

\section*{Definition and Overview}
Angiotensin receptor blockers (ARBs), also known as angiotensin II receptor antagonists or "sartans", are a class of medicines that lower blood pressure by blocking the action of angiotensin II at its receptors. They are used mainly to treat hypertension (high blood pressure) and heart failure, and to protect the kidneys in people with diabetes and chronic kidney disease. Commonly used examples include candesartan, losartan, irbesartan, valsartan, telmisartan, and olmesartan.

ARBs act at a different point in the renin--angiotensin system from ACE inhibitors: rather than reducing the production of angiotensin II, they prevent it from binding to its receptors on blood vessels, the adrenal glands, and the kidneys. They are frequently chosen when a person cannot tolerate an ACE inhibitor, most often because of the dry cough that ACE inhibitors can cause. Their blood pressure lowering, organ-protective, and survival benefits closely match those of ACE inhibitors.

\section*{Epidemiology}
High blood pressure affects around one in three adults worldwide and is a leading modifiable risk factor for stroke, heart attack, heart failure, and chronic kidney disease. ARBs are among the most commonly prescribed blood pressure medicines globally and are increasingly used as first-line treatment, particularly in older adults, in people with diabetes or kidney disease, and in those with heart failure.

Because they rarely cause a cough, ARBs are the usual replacement for people who develop this side effect on an ACE inhibitor. Their use has grown steadily over recent decades, and they now account for a substantial share of all antihypertensive prescribing in many countries.

\section*{Aetiology and Pathophysiology}
The renin--angiotensin--aldosterone system regulates blood pressure and fluid balance. Angiotensin II, produced by the action of angiotensin-converting enzyme, is a powerful vasoconstrictor that also stimulates aldosterone release from the adrenal glands, promoting salt and water retention by the kidney. Overactivity of this system contributes to hypertension, heart failure, and progressive kidney damage.

ARBs interrupt this system at the receptor level. Unlike ACE inhibitors, they do not affect the breakdown of bradykinin, which is why they do not typically cause a dry cough. Their mechanisms include:

\begin{itemize}
    \item \textbf{Blocking the AT1 receptor:} angiotensin II cannot narrow blood vessels, so peripheral resistance and blood pressure fall
    \item \textbf{Reducing aldosterone release:} the kidney excretes more salt and water, reducing fluid overload
    \item \textbf{Lowering glomerular filtration pressure:} pressure within the kidney's filters falls, protecting the kidneys in diabetes and chronic kidney disease
    \item \textbf{Reducing cardiac workload:} afterload falls, and the harmful structural remodelling of the heart in heart failure is slowed
    \item \textbf{Sparing bradykinin:} levels of this vasodilator are not increased, so cough is uncommon
\end{itemize}

\section*{Clinical Features}
Most people taking an ARB have no symptoms, and the benefit is measured through blood pressure, kidney function, and symptoms of heart failure. Possible adverse effects include:

\begin{itemize}
    \item \textbf{Dizziness or light-headedness:} especially on standing, when dehydrated, or after the first dose
    \item \textbf{Hyperkalaemia:} a rise in blood potassium detected on blood tests
    \item \textbf{Rising creatinine:} a small early increase in kidney-function markers that usually stabilises
    \item \textbf{Headache, fatigue, or gastrointestinal upset:} uncommon and usually mild
    \item \textbf{Changes in kidney function:} particularly with dehydration or concurrent kidney-stressing medicines
    \item \textbf{Angioedema:} swelling of the face, lips, or throat; very rare with ARBs but serious if it occurs
    \item \textbf{Dry cough:} distinctly uncommon, which is the main reason ARBs are preferred over ACE inhibitors for some people
\end{itemize}

\section*{Differential Diagnosis}
New symptoms while taking an ARB may be due to the medicine or to another condition:

\begin{itemize}
    \item Dizziness from low blood pressure versus dehydration, anaemia, or a heart rhythm problem
    \item Worsening kidney function from the medicine versus progression of the underlying kidney disease
    \item Raised potassium from the medicine versus renal failure, or an interaction with potassium-sparing diuretics or NSAIDs
    \item Angioedema from the medicine versus allergy to food, insect stings, or another drug
    \item Headache or fatigue from the medicine versus an unrelated illness such as an infection
\end{itemize}

\section*{When to Seek Medical Care}
Tell your doctor if you feel persistently dizzy, faint, or light-headed; if you are started on other medicines that affect the kidneys or potassium; or if you are pregnant, planning a pregnancy, or breastfeeding. Blood tests for kidney function and potassium are needed periodically, and any new persistent symptom should be reviewed.

\textbf{Contact your local emergency services for:}
\begin{itemize}
    \item Swelling of the lips, face, tongue, or throat, or difficulty breathing or swallowing
    \item Fainting, collapse, or severe dizziness
    \item Chest pain, palpitations, or a very fast or irregular heartbeat
\end{itemize}

\section*{Diagnosis and Investigations}
Monitoring follows the same pattern as for ACE inhibitors:

\begin{itemize}
    \item \textbf{Blood pressure measurement:} before starting and at each review to assess the response
    \item \textbf{Kidney function and electrolytes:} creatinine, urea, and potassium are checked before starting and after dose changes
    \item \textbf{Regular monitoring:} kidney function and potassium are rechecked within one to two weeks of starting or increasing the dose, then periodically
    \item \textbf{Pregnancy test:} considered before starting in women of childbearing potential, since ARBs are harmful in pregnancy
    \item \textbf{Medication review:} to detect interactions with NSAIDs, diuretics, potassium supplements, and other blood pressure medicines
\end{itemize}

\section*{Management}
\begin{itemize}
    \item \textbf{Start low and go slow:} the dose is increased gradually to minimise dizziness
    \item \textbf{Consistency:} take the medicine as prescribed, usually once daily
    \item \textbf{Combination therapy:} often combined with a calcium channel blocker or a thiazide diuretic to reach blood pressure targets
    \item \textbf{Switching from an ACE inhibitor:} ARBs are the standard replacement when a dry cough is the problem
    \item \textbf{Heart failure and kidney protection:} ARBs are used long-term at doses optimised by a specialist for these indications
    \item \textbf{Contraindicated in pregnancy:} like ACE inhibitors, ARBs can harm the developing baby and must be stopped before conception
    \item \textbf{Adjunct lifestyle advice:} reduce salt, limit alcohol, stay active, and maintain a healthy weight alongside the medicine
\end{itemize}

\section*{Complications}
\begin{itemize}
    \item Hyperkalaemia, which can cause dangerous heart rhythms in severe cases
    \item Worsening kidney function, especially with dehydration, bilateral renal artery stenosis, or concurrent NSAID use
    \item Hypotension and dizziness, with a risk of falls in older people
    \item Angioedema, which is very rare but can threaten the airway
    \item Fetal harm if used during pregnancy, including effects on the developing kidneys
\end{itemize}

\section*{Prognosis and Outcomes}
ARBs have benefits very similar to those of ACE inhibitors. In heart failure they reduce hospital admissions and improve survival; in diabetes and chronic kidney disease they slow the loss of kidney function and reduce the risk of progression to dialysis; and in hypertension they lower the risk of stroke, heart attack, and heart failure. They are generally very well tolerated, and many people take them for decades with few or no side effects. As with any blood pressure medicine, the benefit depends on taking it regularly and keeping blood pressure in the target range.

\section*{Prevention and Self-Care}
\begin{itemize}
    \item Do not stop the medicine suddenly, since a rebound rise in blood pressure can follow
    \item Keep to scheduled blood tests for kidney function and potassium
    \item Stay well hydrated, but avoid over-the-counter NSAIDs and potassium-rich salt substitutes without medical advice
    \item Rise slowly from sitting or lying to prevent dizziness and falls
    \item Keep an up-to-date medicine list and share it with every health professional you see
    \item If pregnant, planning pregnancy, or breastfeeding, discuss your medicines with your doctor before making changes
    \item Report any swelling of the face, lips, or tongue immediately
\end{itemize}

\section*{Sources and Further Reading}
The following sources provide reliable further information on angiotensin receptor blockers:

\begin{itemize}
    \item NHS. \textit{Information on ARBs, their uses, and possible side effects.} https://www.nhs.uk/conditions/ace-inhibitors/
    \item MedlinePlus. \textit{Overview of angiotensin receptor blockers and their precautions.} https://medlineplus.gov/druginfo/meds/a601092.html
    \item Mayo Clinic. \textit{Information on ARBs and how they compare with ACE inhibitors.} https://www.mayoclinic.org/diseases-conditions/high-blood-pressure/in-depth/ace-inhibitors/art-20048180
    \item NICE. \textit{Clinical guideline on hypertension in adults, including first-line treatment options.} https://www.nice.org.uk/guidance/ng136
    \item MSD Manual. \textit{Angiotensin receptor blockers in the professional health section.} https://www.msdmanuals.com/professional/cardiovascular-disorders/antihypertensive-drugs/angiotensin-ii-receptor-blockers
\end{itemize}
